\section{Introduction}
\label{sec:introduction}

Accurate short-term knowledge of traffic conditions is essential for
understanding and managing the dynamics of urban road networks.
However, traffic evolution is difficult to predict because it results from
the interaction of two strongly coupled components. On the one hand, traffic
at a given location follows temporal patterns determined by its recent
history; on the other hand, changes occurring on one road segment can
propagate through the network and affect other connected locations. A
forecasting model should therefore be able to exploit both temporal
dependencies and the spatial structure of the road network.

Accurate point forecasts alone, however, are not sufficient in practical
settings. Traffic conditions may change rapidly, particularly during
congestion events, making some predictions substantially more uncertain than
others. A forecasting system should therefore provide not only an estimate of
future traffic conditions, but also an indication of the uncertainty
associated with that estimate. Moreover, this uncertainty should remain
meaningful when the model encounters conditions that differ from those
observed during training.

This work investigates these aspects using the METR-LA traffic-sensor
network. Starting from a simple persistence baseline, a temporal forecasting
model based on Gated Recurrent Units (GRUs) is first considered to model the
temporal evolution of each sensor independently. Spatial information is then
introduced through a graph-temporal architecture combining diffusion-style
message passing over the road-sensor graph with recurrent temporal modeling.
Both learned models produce probabilistic forecasts through conditional
quantile estimation, allowing point accuracy and predictive uncertainty to be
evaluated jointly.

Beyond aggregate forecasting performance, particular attention is devoted to
understanding the contribution of the graph and the robustness of the
probabilistic predictions. The graph-temporal model is analyzed at the sensor
level, under different graph configurations and diffusion depths, and under a
distribution shift defined by unusually high network-wide congestion.
Finally, calibration strategies are investigated to determine whether the
reliability of the predictive intervals can be improved when the forecasting
conditions become more challenging.


\subsection{Objectives and Research Questions}
\label{subsec:objectives}

The main objective of this work is to design and evaluate a compact,
probabilistic graph-temporal forecasting pipeline while investigating the
interaction between predictive accuracy, spatial information, uncertainty
calibration, and robustness under distribution shift.

More specifically, the experimental analysis is organized around the
following research questions:

\begin{enumerate}
    \item \textbf{Does graph message passing improve traffic forecasting
    compared with temporal modeling alone?}

    This question is investigated by comparing a persistence baseline, a
    temporal GRU without spatial information exchange, and a graph-temporal
    model that propagates information between connected traffic sensors.

    \item \textbf{How does the graph representation affect forecasting
    performance?}

    The contribution of the graph is examined both at the individual-sensor
    level and through controlled architectural experiments. In particular,
    the analysis investigates where graph message passing helps or hurts,
    whether extending the spatial diffusion depth provides additional
    benefits, and whether the original weighted graph contains useful
    information beyond connectivity alone.

    \item \textbf{How reliable are the probabilistic forecasts under normal
    and shifted traffic conditions?}

    Forecast uncertainty is evaluated through quantile-based predictive
    intervals, considering both empirical coverage and interval width. A
    predefined high-congestion stress test is then used to investigate how
    point accuracy and uncertainty calibration change when the model is
    evaluated under more challenging traffic conditions.

    \item \textbf{Can the calibration degradation observed under
    distribution shift be mitigated?}

    A post-hoc Conformalized Quantile Regression calibration procedure is
    evaluated, while a congestion-aware training strategy is considered as a
    complementary approach for directly improving robustness to rare
    congestion regimes.
\end{enumerate}

Together, these questions shift the analysis beyond average forecasting
accuracy and toward a more complete assessment of when spatial information
is useful, when the model becomes uncertain, and whether its uncertainty
estimates remain reliable under distribution shift.